\documentclass[11pt]{article}
\usepackage{amsmath,amssymb,amsthm,booktabs}
\usepackage[margin=1in]{geometry}

\newtheorem{theorem}{Theorem}
\newcommand{\Psym}{\mathsf P_\infty}
\newcommand{\Mz}{M_\zeta}
\newcommand{\Uinf}{\mathcal U_\infty}

\title{Step 176: Branch C Full-Carrier Foreclosure}
\author{Six Birds Foundations III}
\date{2026-05-15}

\begin{document}
\maketitle

\section{Verdict}

\[
\boxed{\texttt{foreclosure\_verdict = V\_branch\_c\_foreclosed\_generic}.}
\]

The word ``generic'' is used in the diagnostic cascade sense: the obstruction
has now been observed at multiple zeta zeros and for multiple legal Burnol
generators.  It is not an open-dense theorem on the full legal class.

\section{Inherited Numerical Method}

Step 173 supplied the operational identity
\[
K_\infty^{\rm op}(s,s')
=\delta(\tau-\tau')
-\frac{\sin(\lambda(\tau-\tau'))}{\pi(\tau-\tau')}
-\sum_{n\ge0}\Psi_n^\lambda(s)\overline{\Psi_n^\lambda(s')}.
\]
Step 174 gave
\[
L_{\rho,0}(G)
=-\int_{\mathbb R}
\frac{\sin(\gamma-u)}{\pi(\gamma-u)}
\zeta(1/2+iu)G(1/2+iu)\,du
-\sum_{n\ge0}\Psi_n(\rho)J_n
\]
for \(\rho=1/2+i\gamma\), because the delta term cancels at zeta zeros.
Step 175 evaluated the first datum
\[
L_{\rho_1,0}(G_*)
=0.12146653476134234-0.08985731971133760\,i
\]
with error radius \(1.9648066412630187\cdot10^{-4}\) and lower bound
\[
|L_{\rho_1,0}(G_*)|\ge0.1508944091096583.
\]

This step uses the same settings:
\[
\lambda=1,\quad U=200,\quad h=0.05,\quad \texttt{mpmath.zeta precision}=50,
\]
Gauss--Legendre quadrature for \(G\), and sinc-kernel diagonalization on
\([-1,1]\) for the \(\Psi_n\) basis.

\section{Confirmation Data}

The additional legal generator \(G'_*\) uses centers
\[
2,\quad 5/2,\quad 3,
\]
with \(\epsilon=1/4\), and coefficients chosen by the same two-moment system
so that
\[
G'_*(0)=G'_*(1)=0.
\]
The numerical residuals are
\[
G'_*(0)=2.7755575615628914\cdot10^{-17},
\qquad
G'_*(1)=2.7755575615628914\cdot10^{-17}.
\]

\[
\begin{array}{c|c|c|r|r|r}
\toprule
\rho & G & L_{\rho,0}(G) & |L| & \text{error} & \text{lower bound}\\
\midrule
\rho_2 & G_* &
0.16442101662000874-0.027105022872143426\,i &
0.16664019014408440 &
1.9290105611310254\cdot10^{-4} &
0.16644728908797130\\
\rho_3 & G_* &
-0.06557366346002169+0.08988786620051682\,i &
0.11126425225403799 &
1.8585726353641995\cdot10^{-4} &
0.11107839499050157\\
\rho_1 & G'_* &
0.20421356560997495-0.07350308391838643\,i &
0.21703889910486759 &
8.4315870631991870\cdot10^{-5} &
0.21695458323423561\\
\bottomrule
\end{array}
\]

All confirmation values are separated from zero by margins far exceeding their
error radii.

\section{No-Go Theorem}

\begin{theorem}[Branch C full-carrier ZI-COV(i) foreclosure]
On the inherited Burnol/Sonine pulled-evaluator carrier, under the standard
log-Mellin/Fourier normalization for \(\Uinf\) and \(\lambda=1\), the Branch C
full-carrier ZI-COV(i) route is foreclosed.  More precisely, there exist
explicit legal Burnol generators and nontrivial zeta-zero labels
\((\rho,k)=(\rho_j,0)\) for which
\[
L_{\rho,0}(G)
=\langle \Mz G,\Psym y_{\rho,0}\rangle
= (\Psym\Mz G)(\rho)
\ne0.
\]
Consequently, the full-carrier identity
\[
\Psym\Mz G=\Mz A_\infty G
\]
for a bounded full-carrier \(A_\infty\) cannot hold on the full legal Burnol
class.  Step 170's zero-free output subclass remains the valid scope of
ZI-COV(i).
\end{theorem}

\begin{proof}
Step 173 gives the operational identity for \(\Psym\).  Step 174 converts the
zero-evaluator pairing to the computable formula
\[
L_{\rho,0}(G)=-I_{\rm sinc}-\sum_n\Psi_n(\rho)J_n.
\]
Step 175 supplies one certified nonzero datum at \((\rho_1,0,G_*)\).  This
step supplies three further nonzero data:
\[
(\rho_2,0,G_*),\qquad
(\rho_3,0,G_*),\qquad
(\rho_1,0,G'_*).
\]
If full-carrier ZI-COV(i) held, then for every legal \(G\)
\[
\Psym\Mz G=\Mz A_\infty G.
\]
At a zeta zero \(\rho\), the right side has value
\[
(\Mz A_\infty G)(\rho)=\zeta(\rho)(A_\infty G)(\rho)=0
\]
whenever \(A_\infty G\) is a regular legal output at the zero.  The displayed
nonzero values of \(L_{\rho,0}(G)=(\Psym\Mz G)(\rho)\) contradict this
full-carrier route.

Finally, if the normalization is changed by a unitary map \(W\) that is
unitarily equivalent to the standard log-Mellin/Fourier convention, then
\[
\langle f,\Psym y\rangle
=\langle Wf,W\Psym y\rangle
\]
and zero/nonzero status of the matrix element is preserved.  Such a
normalization may alter displayed magnitudes but not the obstruction itself.
\end{proof}

\section{Scope and Nonclaim}

This theorem forecloses only the Branch C full-carrier ZI-COV(i) route.  It
does not close \(\Xi_{\mathrm{BC}}\), does not close Branch A, does not close
Branch B, and does not close any Hecke branch.  It does not prove RH and does
not extend beyond the inherited normalization or normalizations unitarily
equivalent to it.  It also does not decide jet-surjectivity directly; it
decides several matrix-element instances on which that route depends.

\end{document}
